\section{Discusión}
\label{sec:discusion}

\subsection{Hallazgos principales}
\label{sec:dis-hallazgos}
% PRISMA ítem 23a (interpretación general)

Los resultados de la \cref{sec:resultados} se resumen en siete hallazgos. Primero, entre los estudios que abordan la vinculación de esquemas o la autocorrección como contribución central, los multiturno constituyen una fracción minoritaria: D4, proporción descriptiva dentro del marco de elegibilidad, es \DcuatroPct~\% (IC 95~\% de Wilson, \DcuatroICinf--\DcuatroICsup~\%), y la regla registrada lo clasifica como «\DcuatroConclusion{}» (\cref{sec:res-d4}). Segundo, la prueba de tendencia registrada para 2019--2025 no detectó cambios en esa proporción, y la sensibilidad que incluye 2026 apunta en dirección decreciente sin alcanzar significación. Tercero, el historial plano es la representación del estado más frecuente, mientras que las estructuras explícitas y los grafos son minoritarios (\cref{sec:res-pi2}). Cuarto, los contrastes dentro del estudio son mayoritariamente positivos, con medianas de pocos puntos porcentuales (\cref{sec:res-efectos}). Quinto, la inferencia de \en{joins} es implícita en el modelo en la mayoría de los estudios multiturno (\cref{sec:res-pi3}). Sexto, el tipo de historial usado en la evaluación varía entre estudios y con frecuencia no se reporta; solo el estrato con historial predicho permite observar la propagación de errores (\cref{sec:res-pi5}). Séptimo, el número de estudios del corpus que cumplen el criterio pre-especificado de reformulación es \nBrechaMocho{}; este resultado es consistente con prevalencias bajas de cada elemento, porque los pares observados se acercan a los esperados bajo independencia (\cref{sec:dis-brecha}), y no indica que los componentes estén ausentes ni que la combinación se evite.
% fuente: pi15_resultados.json (brecha_M8.componentes), revision_4f.json (A2_M8_esperados, A3_tendencia_2019_2026)

\subsection{Interpretación}
\label{sec:dis-interpretacion}
% PRISMA ítems 23a (interpretación) y 23b (limitaciones de la evidencia)

La conclusión de que el escenario multiturno está «\DcuatroConclusion{}» procede de la aplicación de una regla registrada a un intervalo de confianza, no de una prueba de hipótesis. Se trata, además, de una referencia condicionada al marco de elegibilidad: los estudios de un turno debían cumplir más requisitos que los multiturno, lo que tiende a aumentar la proporción multiturno. Por ello, el indicador no se generaliza a todo el campo del Text-to-SQL. Las variantes de sensibilidad con conclusión registrada conservan la clasificación. La cota extrema que incorpora los informes no recuperados también queda por debajo del umbral, con un límite superior de 19,1~\%, aunque para esa variante no se fijó una regla de antemano; la estimación probabilística, que aplica las tasas observadas entre los informes recuperados, da D4 $\approx$ \DcuatroProbPct~\%. La ausencia de tendencia detectada ($p$ = \DcuatroTendenciaP) coexiste con intervalos anuales amplios: entre 2019 y 2024 el límite superior contiene el \umbralDcuatro{}~\%, de modo que la potencia de la prueba en los primeros años es baja. La sensibilidad que incluye 2026 da $p$ = \tendenciaConP{}, en dirección decreciente y no significativa al 0,05; se conserva la prueba registrada como principal. Los datos no respaldan, por tanto, que la proporción multiturno esté aumentando, pero tampoco permiten afirmar que disminuya, porque 2026 es un año parcial.
% fuente: pi0_resultados.json, sensibilidad.json, PI0.md, revision_4f.json (A3_tendencia_2019_2026, A7_cota_probabilistica)

Las técnicas de un turno aparecen en multiturno con frecuencias relativas distintas (\cref{tab:contraste-pi34}, con ambos brazos restringidos a estudios con PI3/PI4 centrales: \nCentralMT{} multiturno y \nCentralUT{} de un turno). La aclaración con el usuario y la señal de usuario siguen siendo proporcionalmente mucho más frecuentes en multiturno (47,1 y 44,1~\%, frente a 2,2 y 3,0~\%), y el verificador o crítico es mucho menos frecuente (2,9~\% frente a 21,6~\%). La retroalimentación de ejecución está más cerca entre brazos que en el contraste previo sin restringir (23,5~\% frente a 33,6~\%). La autoconsistencia por votación y la reparación simbólica no figuran en ningún estudio multiturno de la tabla, si bien 2 estudios multiturno evalúan componentes de selección de candidatos (\cref{tab:efectos}). La inferencia de \en{joins}, en cambio, no es comparable: la fila «no reportado o no extraído» es mucho mayor en un turno (72,0~\% frente a 17,6~\%) porque el Nivel 1 se extrajo de las páginas iniciales, de modo que la mayor frecuencia de inferencia implícita en multiturno puede ser un efecto de la extracción. Se trata de frecuencias de uso y no de resultados, y la base del contraste difiere de la base con umbral de validez de las taxonomías. En cuanto a la inferencia algorítmica de \en{joins}, es casi inexistente en multiturno: solo 2 estudios usan una inferencia no implícita~\cite{Dou2023UniSAr,Wang2020PGGSQL}. Con 2 estudios no es posible contrastar, dentro de un mismo estudio, la inferencia implícita con la algorítmica, y los archivos no incluyen un análisis por tamaño del esquema. Por consiguiente, la evidencia no permite valorar si la inferencia implícita es suficiente en esquemas grandes.
% fuente: pi34_tecnicas_central.csv, PI3.md, PI4.md

Los componentes de grafo o estado estructurado muestran dirección positiva en los 11 estudios que los evalúan, con confianza registrada alta, que baja a moderada al penalizar el sesgo de notificación (\cref{tab:efectos}). Esta evidencia se acota con tres salvedades. La primera es el sesgo de notificación: las ablaciones suelen reportarse para justificar el componente propuesto, de modo que los contrastes negativos son escasos; con el umbral de nulidad de $|\Delta| < 0{,}5$~pp, \nulosUmbral{} de los \nEstClase{} resultados estudio-clase son nulos, y la confianza registrada es «alta» en \nConfAlta{} clases pero baja a «moderada» al rebajarla un nivel (\cref{sec:res-sesgos}). La segunda es que algunas medianas cambian al excluir los \en{preprints}; la de SQL o resultados previos, por ejemplo, pasa de 5,7 a 3,3~pp. La tercera es que la clase «otros componentes» es heterogénea. Lo que la evidencia muestra es que, cuando se reporta una ablación, el componente estructurado mejora el resultado frente a su propia variante sin él. No muestra que el estado estructurado sea la causa de la mejora en general, ni que supere al historial plano en comparaciones entre estudios.
% fuente: pi_efectos_componentes.csv, confianza_hallazgos.csv, sensibilidad.json, revision_4f.json (A4_nulos, A4_confianza_penalizada)

Una parte considerable de la información relevante figura como no reportada o no extraída: la técnica de restauración del contexto, el mecanismo de autocorrección, el historial de la evaluación y los campos descriptivos del grafo (\cref{sec:res-sesgos}). Esta ausencia limita las taxonomías y la comparabilidad. Sin embargo, una fracción de los valores no reportados podría deberse a la extracción asistida por IA y no a los informes. El acuerdo entre extractores en los campos codificados del Nivel 2 fue del \acuerdoNivelDos~\%, y en la muestra independiente de la autora el historial coincidió en 3 de 5 estudios; la estratificación de PI5 depende de esa codificación. Respecto de la evaluación, el historial de referencia no expone la propagación de errores, porque el sistema nunca recibe sus propias consultas erróneas. Se plantea, como hipótesis, que la evaluación con historial \en{gold} sobreestime la robustez. El corpus no permite contrastarla: ninguna curva de degradación con al menos tres turnos corresponde a ese estrato. Más aún, el estrato predicho reúne solo 6 curvas de 5 estudios y se degrada menos (mediana de $-7{,}0$~pp por turno, rango de \degPredMin{} a \degPredMax{}) que el de solo preguntas (mediana de $-11{,}0$~pp, rango de \degSoloMin{} a \degSoloMax{}), lo contrario de lo esperable si la propagación de errores fuera el factor dominante. Con tan pocas curvas, \en{benchmarks} y métricas mezclados y una profundidad de turno confundida con la dificultad de las preguntas, esta diferencia no permite concluir y puede deberse a esos factores.
% fuente: informe_acuerdo_extraccion.md, PI5.md, revision_4f.json (A9_degradacion)

Por último, el análisis bibliométrico exploratorio indica que los estudios multiturno se publican sobre todo en foros de procesamiento del lenguaje natural (34) y en mucha menor medida en foros de bases de datos (4), con 18 \en{preprints} (\cref{fig:bib_panorama}, panel~(a)). Este reparto es coherente con que el problema se formule principalmente como un problema de lenguaje, y coexiste con la escasa inferencia algorítmica de \en{joins} observada en multiturno, sin que se haya evaluado relación alguna entre ambos hechos. El análisis es exploratorio (enmienda A6) y depende de la clasificación de los foros y de la cobertura de OpenAlex.
% fuente: bib_comunidad.csv

\subsection{Comparación con revisiones previas}
\label{sec:dis-comparacion}
% PRISMA ítem 23a (interpretación en el contexto de otra evidencia)

La comparación con las revisiones de la \cref{tab:relacionados} es de alcance y de método, ya que sus objetivos difieren de los de esta revisión. Según su resumen, Yang et al.~\cite{Yang2025Conversational} cubren el manejo del contexto, la evolución de las arquitecturas y una comparación cuantitativa en SParC y CoSQL, sin mencionar protocolo de búsqueda ni tratar como ejes la vinculación de esquemas, la inferencia de \en{joins} o la retroalimentación de ejecución. Frente a esa comparación de resultados, la presente revisión aporta recuentos sistemáticos y contrastes dentro del estudio, y evita comparar resultados entre estudios. No se puede afirmar qué cubre su texto completo, que no fue posible leer. Li, Qu y Haffari~\cite{Li2020ContextSurvey} ofrecen una visión narrativa anterior a los LLM y más amplia que SQL, mientras que el corpus actual abarca la transición hacia el \en{prompting} y los agentes (\cref{fig:pi0_paradigma_anio}). Entre las revisiones sistemáticas de 2026, la versión arXiv v2 de Liu et al.~\cite{Liu2026NLIDBReview} no documenta protocolo ni menciona SParC ni CoSQL. Bui et al.~\cite{Bui2026Survey} documentan una búsqueda sistemática, pero en el manuscrito leído el escenario multiturno aparece solo como categoría de \en{datasets}. Ortiz-Delgado et al.~\cite{OrtizDelgado2026RAG} no tratan el escenario multiturno, según su resumen y su registro.

Los aportes diferenciales son, por tanto, el registro previo del protocolo, el alcance definido con PICOC, el contraste explícito entre un turno y multiturno, la degradación estratificada por historial y los contrastes dentro del estudio con un nivel de confianza por hallazgo. Ello no implica superioridad sobre los \en{surveys}, cuyas taxonomías de arquitecturas y panoramas conceptuales cumplen otra función. La contribución de las fuentes de búsqueda también aporta información de método (\cref{fig:busqueda_fuentes}). IEEE Xplore aportó 18 estudios incluidos a partir de \nBaseIEEE{} registros, lo que es compatible con el sesgo de disponibilidad declarado en la enmienda A4: los informes de editoriales de pago están infrarrepresentados entre los textos completos recuperados. Con los datos disponibles, sin embargo, ese efecto no puede separarse de una menor pertinencia de los registros. arXiv, en cambio, aportó 244 estudios incluidos, 125 de ellos recuperados solo por esa fuente, lo que respalda la decisión de incluir \en{preprints}. Bui et al.~\cite{Bui2026Survey} también reportan arXiv como la fuente con más estudios. El alcance de estos sesgos se analiza en la \cref{sec:amenazas}.
% fuente: busqueda_fuentes.csv, enmiendas.md (A4), 00_analisis_brecha.md

\subsection{Criterio de reformulación y brechas de cobertura}
\label{sec:dis-brecha}
% PRISMA ítem 23d (implicaciones para la investigación futura)

El criterio pre-especificado de reformulación evaluaba la novedad de una conjunción. Se registró en el protocolo (v1.2, M8a), con las definiciones de los elementos (a), (b) y (c), como aquello que contaría en contra de la brecha planteada de antemano, y se reporta con independencia del resultado. Combina tres elementos: (a) un estado estructurado explícito, (b) la inferencia algorítmica de \en{joins} o la autocorrección por ejecución, y (c) la evaluación con historial predicho o solo con preguntas. El criterio se habría activado con al menos 2 estudios que reunieran los tres elementos. Entre los \nMultiturno{} estudios multiturno del Nivel 2, \nMochoA{} presentan (a), \nMochoB{} presentan (b) y \nMochoC{} presentan (c). Los pares observados son \nMochoAB{} para (a) y (b) (esperado bajo independencia, \espMochoAB{}), \nMochoAC{} para (a) y (c) (esperado, \espMochoAC{}) y \nMochoBC{} para (b) y (c) (esperado, \espMochoBC{}). El número de estudios que reúnen los tres elementos es \nBrechaMocho{} (esperado bajo independencia, \espMochoABC{}), y la sensibilidad prevista con grafos neuronales tampoco identificó ninguno. Los pares observados se acercan, por tanto, a lo que predice la independencia: los datos muestran bajas prevalencias de cada elemento, no una combinación evitada. Cada componente existe en el corpus, y los componentes figuran asimismo por separado en la literatura: representaciones de la interacción en forma de grafo~\cite{Cai2020IGSQL}, que como codificador neuronal no cuentan para (a); inferencia no implícita de \en{joins} en multiturno~\cite{Dou2023UniSAr,Wang2020PGGSQL}; y rutas de \en{join} algorítmicas en un turno~\cite{Safdarian2026SchemaGraphSQL}. Un resultado nulo sobre una conjunción es evidencia débil: cuanto más específica es la combinación de rasgos exigida, menos probable es que un estudio la reúna, con independencia de su valor. El resultado indica que la combinación no se encontró en el corpus; no indica que sea necesaria ni que vaya a resultar eficaz. Es, además, condicional por tres razones. En primer lugar, depende de la codificación de la extracción: un estudio que combine los elementos sin reportarlos con claridad, o cuyos campos se codificaran como no reportados o no extraídos, no se detectaría. En segundo lugar, los informes no recuperados no se evaluaron, y 43 de ellos tenían señal multiturno (asignada por título y resumen). En tercer lugar, el acuerdo entre extractores en el Nivel 2 fue moderado. Es, por tanto, una afirmación sobre el corpus de esta revisión y no sobre toda la literatura.
% fuente: pi15_resultados.json (brecha_M8, brecha_M8.componentes), revision_4f.json (A2_M8_esperados), PI2.md, pi0_resultados.json, 01_preguntas_picoc.md

A partir de la evidencia se identifican las siguientes brechas de cobertura observadas. Se ordenan de mayor a menor solidez según una regla: primero las sustentadas por datos que los informes declaran directamente, luego las que dependen de la codificación de la extracción y, al final, las basadas en la rareza. Se presentan como ausencias o escaseces en el corpus, no como necesidades ni como soluciones cuya eficacia esta revisión haya establecido, y para cada una se indica qué muestra la evidencia y qué no permite mostrar:
\begin{enumerate}
  \item Evaluación con historial predicho y reporte por turno. El historial de la evaluación no se declara en 20 de 72 estudios, y el estrato predicho reúne solo 6 curvas de degradación de 5 estudios (\cref{sec:res-pi5}). La evidencia procede directamente de lo que los informes declaran; no permite estimar cuánto sobreestima la robustez la evaluación con historial de referencia.
  % fuente: pi5_historial.csv, pi15_resultados.json, PI5.md
  \item Reporte de la técnica y de la evaluación. Solo 14 de los estudios del Nivel 2 (18,2~\%) describen la técnica con detalle suficiente para reproducirla (CQ2), y solo 19 (24,7~\%) reportan a la vez métricas estándar y el tipo de historial (CQ5); varios campos de la extracción quedan como no reportados o no extraídos (\cref{sec:res-sesgos}). La evidencia muestra limitaciones de reporte; no permite separar del todo las omisiones de los informes de las atribuibles a la extracción asistida por IA.
  % fuente: CALIDAD.md
  \item Contaminación de los \en{benchmarks}. Solo 7 de los estudios que alcanzan el umbral de validez (9,7~\%) la discuten. La evidencia muestra que el riesgo rara vez se examina; no permite estimar cuánto infla los resultados reportados (\cref{sec:am-externa}).
  % fuente: pi5_historial.csv
  \item Datos en español y multilingües. Solo \nDatasetES{} de los \nEstudios{} estudios incluidos evalúan con datos en español, frente a \nDatasetEN{} con datos en inglés y \nDatasetZH{} en chino. La evidencia muestra una cobertura lingüística desigual, en un corpus que además excluye publicaciones en otros idiomas; no muestra si el rendimiento o la dificultad del diálogo difieren entre idiomas. Como las cadenas de búsqueda estaban en inglés, las publicaciones en español pueden haberse omitido; los conjuntos de datos en español descritos en artículos en inglés sí se habrían encontrado.
  \item La aclaración con el usuario como ruta alternativa. Figura en el 47,1~\% de los estudios multiturno frente al 2,2~\% de los de un turno, y la señal de usuario en el 44,1~\% frente al 3,0~\% (\cref{tab:contraste-pi34}). La evidencia muestra que la interacción es una vía frecuente para resolver la ambigüedad en multiturno; no muestra su eficacia relativa, porque los contrastes dentro del estudio comparan cada componente con su ausencia y no con otras rutas, como la autocorrección por ejecución o el estado estructurado.
  % fuente: pi34_tecnicas_central.csv
  \item Estado estructurado combinado con inferencia algorítmica de \en{joins} o retroalimentación de ejecución. El número de estudios multiturno del Nivel 2 que combinan (a) y (b) es \nMochoAB{} (esperado bajo independencia, \espMochoAB{}); por separado, la inferencia no implícita de \en{joins} aparece en 2 estudios multiturno y la retroalimentación de ejecución en el 23,5~\% de los multiturno con PI3/PI4 centrales, frente al 33,6~\% de los de un turno. Se trata de una rareza, no de una necesidad establecida: la revisión no permite valorar si la inferencia implícita de \en{joins} es insuficiente ni si la combinación mejora los resultados (\cref{sec:dis-interpretacion}).
  % fuente: pi15_resultados.json (brecha_M8.componentes), pi34_tecnicas_central.csv, PI3.md, revision_4f.json (A2_M8_esperados)
  \item Trabajo sin entrenamiento. Su presencia aumenta entre los estudios con vinculación de esquemas o autocorrección central (\cref{sec:res-caracteristicas}), pero los archivos no la desagregan por escenario, de modo que su alcance en multiturno no está caracterizado.
\end{enumerate}

\subsection{Implicaciones para el reporte y la evaluación}
\label{sec:dis-practica}
% PRISMA ítem 23d (implicaciones para la práctica)

De los resultados se derivan recomendaciones de reporte para los estudios primarios. Conviene declarar siempre el tipo de historial usado en la evaluación y reportar los resultados por turno, preferentemente con historial predicho, que es el único que refleja la propagación de errores. Junto con la exactitud de ejecución, conviene reportar las métricas por pregunta y por interacción, que capturan la coherencia a lo largo del diálogo. También deberían describirse la técnica de restauración del contexto, la representación del estado y su actualización, el coste del contexto en función del número de turnos, y el mecanismo y la señal de autocorrección; su omisión contribuye a los valores no reportados de esta revisión.

En cuanto a la reproducibilidad, entre los estudios del Nivel 2 la técnica se describe con detalle suficiente para reproducirla en solo 14 (18,2~\%), y el código o los datos no están disponibles en 24 (31,2~\%) (\cref{sec:res-calidad}). Reportar las ablaciones con el mismo modelo base, incluidas las de efecto nulo o negativo, atenuaría el sesgo de notificación señalado en la \cref{sec:dis-interpretacion}. Solo 7 de los estudios que alcanzan el umbral de validez (9,7~\%) discuten la contaminación del \en{benchmark}, lo que impide valorar su influencia en los resultados reportados. Por último, para quienes diseñan \en{benchmarks} y protocolos de evaluación, agrupar los turnos con una dificultad controlada permitiría separar el efecto de la profundidad del turno del de la dificultad de las preguntas, una distinción que las pendientes actuales no permiten hacer.
% fuente: CALIDAD.md, PI5.md
